%! TeX program = lualatex
% The three cases of the first derivative test, each drawn as the number line
% the lesson tells students to draw: the critical number in the middle, the
% sign of f' on either side, and an arrow for what that sign does to f.
%
%   page 1 -- positive then negative, a local maximum
%   page 2 -- negative then positive, a local minimum
%   page 3 -- positive twice, neither
%   page 4 -- negative twice, neither
%
% The third case needs two pages because "the same sign" is two pictures, not
% one: the curve can rise through c or fall through c, and in both the
% tangent at c is horizontal without an extremum there.  Showing only one of
% them invites a reader to think the case is about decreasing functions.
%
% The line carries a single arrow, on the right, labelled x -- the same way
% pgfplots draws an x-axis everywhere else in these notes.  Everything is
% \footnotesize so a pair of them fits side by side in the text column.
%
% [tikz] rather than [multi=tikzpicture]: the latter does not load TikZ
% itself, and this file needs \usetikzlibrary before the preamble input.
\documentclass[tikz, crop]{standalone}
\usetikzlibrary{arrows.meta}

\input{../typography.tex.preamble}

% The axis itself: one arrow to the right, a tick at the critical number, c
% beneath it.  Not called \numberline -- LaTeX already owns that name for
% table of contents entries.
\newcommand{\signaxis}{%
  \draw[-{Stealth}] (-1.6,0) -- (1.6,0) node[right] {\(x\)};
  \draw[thick] (0,-0.1) -- (0,0.1);
  \node[below] at (0,-0.12) {\(c\)};}

% One side: colour, x position, the relation symbol, and the slope arrow that
% the sign produces.  Sign above the line, slope below, so a reader can take
% the two rows separately.
\newcommand{\sidelabel}[4]{%
  \node[#1, above] at (#2,0.08) {\(f' #3 0\)};
  \node[#1, below] at (#2,-0.12) {\(#4\)};}

\begin{document}

% ---------------------------------------------------------------------
% Case 1 -- f' changes from positive to negative: a local maximum.
\begin{tikzpicture}[font=\footnotesize]
  \signaxis{}
  \sidelabel{main}{-0.8}{>}{\nearrow}
  \sidelabel{supp}{0.8}{<}{\searrow}
\end{tikzpicture}

% ---------------------------------------------------------------------
% Case 2 -- f' changes from negative to positive: a local minimum.
\begin{tikzpicture}[font=\footnotesize]
  \signaxis{}
  \sidelabel{supp}{-0.8}{<}{\searrow}
  \sidelabel{main}{0.8}{>}{\nearrow}
\end{tikzpicture}

% ---------------------------------------------------------------------
% Case 3, first possibility -- positive on both sides: f rises through c.
\begin{tikzpicture}[font=\footnotesize]
  \signaxis{}
  \sidelabel{main}{-0.8}{>}{\nearrow}
  \sidelabel{main}{0.8}{>}{\nearrow}
\end{tikzpicture}

% ---------------------------------------------------------------------
% Case 3, second possibility -- negative on both sides: f falls through c.
\begin{tikzpicture}[font=\footnotesize]
  \signaxis{}
  \sidelabel{supp}{-0.8}{<}{\searrow}
  \sidelabel{supp}{0.8}{<}{\searrow}
\end{tikzpicture}
\end{document}
